\begin{theorem}[MetaCIC parallel-diamond frontier seed closure marker]
\label{thm:metacic-parallel-diamond-frontier-seed-closure-marker}
Every seed registration route for $\MetaCICParallelDiamondFrontierUp$ reads the carrier row, critical-pair envelope, candidate-normalization scope, residual-substitution dependency, finite residual witness, local join, obstruction, transport, replay, provenance, and local $\NameCert$ rows before any candidate-mediated strong-normalization or confluence consumer may use the frontier. In the displayed spine
$$
P,K,J,R,S,C,B,O,H,T,G,N,
$$
$P$ retains the residual-substitution dependency, $K$ exposes the critical-pair envelope, $J$ records the candidate-normalization row, $R,S,C,B$ carry the finite residual witness and local-join budget, $O$ retains the obstruction row, and $H,T,G,N$ provide transport, replay, provenance, and local naming.
\end{theorem}
\begin{proof}
Start with the frontier carrier and obligation surface. They display the finite rows $P,K,J,R,S,C,B,O,H,T,G,N$ and keep the residual-substitution premise, critical-pair data, candidate row, residual frontier, checker rows, obstruction, transport, replay, provenance, and local name as ordinary packet coordinates.

Read the seed-facing child routes next. The critical-pair envelope, candidate-normalization scope, residual-substitution dependency route, finite residual witness consumption, and residual-join locality rows force the seed route to pass through the same displayed packet before a consumer sees the frontier.

Finally apply the closurestatus guard and candidate-SN handoff already present in the chapter. They keep the obstruction row and structural coordinates visible before any candidate-mediated strong-normalization or confluence consumer can use the route. Hence the seed closure marker remains a finite frontier registration surface, not residual-substitution discharge, subject reduction, full Church--Rosser, evaluator equality, quotient equality, or kernel consistency.
\end{proof}
